\section{Related Evidence and the Empirical Object}\label{sec:literature}

The closest motivating study is the August 2026 revision of \citet{brynjolfsson2026canaries}.  Using ADP payroll records through June 2026, it reports weaker employment growth among ages 22--25 in occupations with high GPT-4 beta exposure and studies employer-match hiring and separation.  The CPS cannot reproduce its employer panel, standardized-title mapping, firm controls, or worker--firm flows.  Because that study also reports public CPS and ACS evidence, novelty here lies in auditing a fixed national-CPS estimand, not discovering the age pattern.

Administrative and workplace studies observe adoption or tool access in particular settings \citep{humlum2025still,brynjolfsson2025generative}.  Firm-level evidence also distinguishes direct task substitution, reallocation across tasks, and firm-growth effects \citep{hampole2025artificial}.  Survey evidence shows that realized generative-AI use varies substantially even among workers doing similar work and that exposure explains only part of adoption \citep{bick2026work}.  Occupation scores instead describe tasks a technology might affect.  Remote-work evidence supplies another competing interpretation \citep{emanuel2026proximity}.  None turns a capability score into observed deployment.

The primary score is Eloundou et al.'s GPT-4 beta: tasks are labeled by whether an LLM, alone or with complementary software, could reduce completion time at constant quality \citep{eloundou2024gpts}.  AIOE maps AI applications to abilities \citep{felten2018method,felten2021occupational}; Webb measures patent--task overlap \citep{webb2020impact}; Dingel--Neiman classifies teleworkability \citep{dingel2020home}.  Computer use, routine content, education, and remotability are therefore overlapping occupational characteristics, not clean alternative treatments.  The task literature likewise allows technology to substitute for some tasks, complement others, and reallocate work within occupations \citep{autor2003skill,autor2013growth,acemoglu2019automation}.  Faster-changing technical careers may also combine early entry advantages with skill obsolescence later in the career \citep{deming2020earnings}, while the employment consequence of automating a task depends on whether the remaining bundle requires more or less expertise \citep{autor2025expertise}.  These mechanisms make an ages-22--25 versus older-worker contrast economically relevant but not uniquely diagnostic of entry-level AI substitution.

This distinction changes the burden of interpretation.  Agreement across correlated scores can show that a direction is not unique to one coding rule, but cannot validate a latent exposure construct.  Conversely, coefficient differences can arise from task definitions, occupational support, or treatment-group membership.  The empirical work below therefore matches supports where possible and names population changes where it is not.

I distinguish six layers:
\[
C\longrightarrow M\longrightarrow H\longrightarrow \Omega\longrightarrow R\longrightarrow Q,
\]
the intended construct, measurement rule, occupational harmonization, retained support, representation, and regression estimand.  A vintage mismatch is an implementation error; score repair on fixed support changes measurement; occupation readmission changes population; and a different exposure contrast changes the estimand.  Common-draw paired inference tests coefficient movement, while each component interval measures its own resolution.  Neither a detected movement nor a nondetection assigns a causal share or establishes equivalence.  This focus connects the exercise to specification-sensitivity and generated-covariate inference \citep{andrews2017sensitivity,rambachan2023credible,battaglia2024inference,christensen2026valid,duan2026inference,ludwig2026large}.
